% Source: https://artofproblemsolving.com/community/c3h2096167_number_theory_fun
% Stable identity: aops-topic-2096167-q04
We have $1001=7\cdot11\cdot13$ and $n^3-n=(n-1)n(n+1)$. For a prime $p$, $p$ divides $(n-1)n(n+1)$ exactly when it divides one of the factors, that is, when $n\equiv0,1$ or $-1\pmod p$: $3$ residues modulo each of $7$, $11$ and $13$.

By the Chinese remainder theorem, the conditions modulo $7$, $11$ and $13$ combine into exactly $3\cdot3\cdot3=27$ residue classes modulo $1001$. The numbers $1,2,\dots,1000$ contain each residue class modulo $1001$ exactly once, except the class of $0$. The class of $0$ is one of the $27$ (as $0\equiv0$ modulo each prime), so the count is $27-1=\boxed{26}$.
